\documentclass[12pt, twoside]{article}
\input{../preamble}

\fancyhead[LE]{\thepage}
\fancyhead[RO]{Markscheme}
\fancyhead[LO]{La Scuola d'Italia / Dr.\ Huson / IB Math \\* 17 April 2026}

\begin{document}
\subsubsection*{6.8 Quiz: Polynomials Challenge --- Markscheme}

\begin{enumerate}[itemsep=0.3cm]

%--- Problem 1: f(x) = -2(x+2)(x-4) [7 marks] ---
%--- Modeled after AA SL 2021 Nov P1 Q01 ---
\item[1.] [Maximum mark: 7]

Consider the function $f(x) = -2(x + 2)(x - 4)$, for $x \in \mathbb{R}$.

\medskip

\textbf{(a)(i)} Write down the $x$-intercepts. \hfill [2]

$x = -2$ and $x = 4$ \hfill \textbf{A1A1}

\medskip

\textbf{(a)(ii)} Find the coordinates of the vertex. \hfill [3]

Axis of symmetry: $x = \dfrac{-2 + 4}{2} = 1$ \hfill \textbf{M1}

\emph{M1 may also be awarded for any attempt at completing the square.}

$f(1) = -2(1 + 2)(1 - 4) = -2(3)(-3) = 18$

Vertex: $(1, 18)$ \hfill \textbf{A1A1}

\medskip

\textbf{(b)} $f(x) = -2(x - h)^2 + k$. Write down $h$ and $k$. \hfill [2]

$h = 1$ \hfill \textbf{A1}

$k = 18$ \hfill \textbf{A1}

\emph{FT: award A1A1 for $h$ and $k$ read from the candidate's vertex in (a)(ii).}

\bigskip
\hrule
\bigskip

%--- Problem 2: y = mx tangent to y = x^2 + 9 [7 marks] ---
\item[2.] [Maximum mark: 7]

The line $y = mx$ is tangent to the parabola $y = x^2 + 9$.

\medskip

\textbf{(a)} Show that at any intersection, $x^2 - mx + 9 = 0$. \hfill [1]

$x^2 + 9 = mx$ \hfill \textbf{M1}

$x^2 - mx + 9 = 0$ \hfill \textbf{AG}

\medskip

\textbf{(b)} Find the discriminant in terms of $m$. \hfill [1]

$\Delta = (-m)^2 - 4(1)(9) = m^2 - 36$ \hfill \textbf{A1}

\medskip

\textbf{(c)} Find the values of $m$ for tangency. \hfill [2]

Tangent: $\Delta = 0$ \hfill \textbf{M1}

$m^2 - 36 = 0$

$m = \pm 6$ \hfill \textbf{A1}

\medskip

\textbf{(d)} Find the coordinates of each point of tangency. \hfill [3]

For $m = 6$: $x^2 - 6x + 9 = 0$

$(x - 3)^2 = 0$ \hfill \textbf{M1}

$x = 3$

$y = 6(3) = 18$. \quad Point of tangency: $(3, 18)$ \hfill \textbf{A1}

For $m = -6$: $x^2 + 6x + 9 = 0$

$(x + 3)^2 = 0$

$x = -3$

$y = -6(-3) = 18$. \quad Point of tangency: $(-3, 18)$ \hfill \textbf{A1}

\emph{FT: award A1 for each point correctly computed from the candidate's $x$-value.}

\newpage

%--- Problem 3: f(x) = 2x^2 - 8x + 3 [7 marks] ---
\item[3.] [Maximum mark: 7]

Consider the function $f(x) = 2x^2 - 8x + 3$.

\medskip

\textbf{(a)} Write $f(x)$ in the form $a(x - h)^2 + k$. \hfill [3]

$f(x) = 2(x^2 - 4x) + 3$ \hfill \textbf{M1}

\emph{M1 for any attempt at completing the square, including use of $-b/(2a)$ anywhere.}

$= 2(x^2 - 4x + 4 - 4) + 3 = 2(x - 2)^2 - 8 + 3$

$= 2(x - 2)^2 - 5$ \hfill \textbf{A1A1}

\medskip

\textbf{(b)} Write down the coordinates of the vertex. \hfill [1]

Vertex: $(2, -5)$ \hfill \textbf{A1}

\medskip

\textbf{(c)} Find the discriminant of $f(x) = 0$. \hfill [1]

$\Delta = (-8)^2 - 4(2)(3) = 64 - 24 = 40$ \hfill \textbf{A1}

\medskip

\textbf{(d)} State the number of $x$-intercepts and justify. \hfill [2]

Two $x$-intercepts. \hfill \textbf{A1}

\emph{Accept ``two'', ``2'', ``double'', or equivalent.}

Justification: $\Delta = 40 > 0$, so two distinct real roots. \hfill \textbf{R1}

\emph{R1 may also be awarded for naming the two roots, or for stating that the candidate's $\Delta > 0$.}

\bigskip
\hrule
\bigskip

%--- Problem 4: p(x) = 2x^3 + x^2 - 13x + 6 [8 marks] ---
\item[4.] [Maximum mark: 8]

Let $p(x) = 2x^3 + x^2 - 13x + 6$.

\medskip

\textbf{(a)} Show that $(x + 3)$ is a factor of $p(x)$. \hfill [2]

$p(-3) = 2(-27) + (-3)^2 - 13(-3) + 6$ \hfill \textbf{M1}

$= -54 + 9 + 39 + 6 = 0$

Since $p(-3) = 0$, $(x + 3)$ is a factor by the
factor theorem. \hfill \textbf{A1 AG}

\medskip

\textbf{(b)} Express $p(x) = (x + 3)(ax^2 + bx + c)$. \hfill [2]

By polynomial division: \hfill \textbf{M1}

$2x^3 + x^2 - 13x + 6 = (x + 3)(2x^2 - 5x + 2)$ \hfill \textbf{A1}

\emph{Accept $a = 2$, $b = -5$, $c = 2$.}

\emph{M1 may also be awarded for synthetic division (Ruffini's rule).}

\medskip

\textbf{(c)} Factorise $p(x)$ completely. \hfill [2]

$2x^2 - 5x + 2 = (2x - 1)(x - 2)$ \hfill \textbf{M1}

\emph{M1 may also be awarded for use of the quadratic formula.}

$p(x) = (x + 3)(2x - 1)(x - 2)$ \hfill \textbf{A1}

\emph{Accept $2(x - \tfrac{1}{2})$ in place of $(2x - 1)$; do not accept $(x - \tfrac{1}{2})$ alone (missing leading coefficient).}

\medskip

\textbf{(d)} Write down all solutions to $p(x) = 0$. \hfill [1]

$x = -3, \quad x = \dfrac{1}{2}, \quad x = 2$ \hfill \textbf{A1}

\medskip

\textbf{(e)} Write down the $y$-intercept. \hfill [1]

$p(0) = 6$, so the $y$-intercept is $(0, 6)$. \hfill \textbf{A1}

\end{enumerate}
\end{document}
